\section{Related Work}

\paragraph{E-graphs and Equality Saturation}
\Egraphs were first introduced by \citet{nelson}
 in late 1970s 
 to support a decision procedure for the theory of equalities.
\citet{tarjan-congruence} later introduced
 a more efficient algorithm
 and analyzed the its time complexity.
\Egraphs are used at the core of 
 many theorem provers and solvers~\citep{simplify, z3, cvc4}.
Because \egraphs can compactly represent program spaces,
 they were repurposed for program optimization in the 2000s~\citep{eqsat,denali}.
Other data structures for compact program space representations are developed in parallel,
 including finite tree automata~\citep{dace,blaze} and version space algebras~\citep{vsa,flashmeta}.
There are two essential problems to these data structures:
 how to construct the desired program space and how to search it.
\citet{eqsat} observed that 
 the program space can be grown via equational, non-destructive rewrites, 
 which they called equality saturation.
This insight
 leads to a line of work on using equality saturation 
 for program optimization and program synthesis~\citep{herbie,ruler,tensat,spores,diospyros,szalinski}.
However, 
 a problem with this rewriting-based approach to program space construction is that,
 in many cases,
 sound rewrite rules are difficult to define in a purely syntactic way.
The \egg framework by \citet{egg} mitigates this issue 
 by introducing \eclass analyses, 
 which allow simple semantic analyses over the \egraphs.
Our work improves \eclass analyses 
 by allowing more expressive analysis rules to be defined
 compositionally.

\citet{relational-ematching} first studied 
 the connection between \egraphs and relational databases.
By reducing pattern matching over \egraphs to relational queries,
 they made the matching procedure orders of magnitude faster.
However, their technique has the dual representation problem, 
 i.e., one has to keep both the \egraph and its relational representation, 
 which limits its practical adoptions.
We build on their work
 and view the entire equality saturation algorithm from the relational perspective.
This saves us from synchronizing two representations of \egraphs
 and further exploits the performance benefits of the relational approach.

\Egglog also brings new insights on some problems in \egraphs and equality saturation.
For example, 
 the literature studied the problem of incremental pattern matching over \egraphs.
\citet{relational-ematching} conjectured that this problem can be solved by 
 classical techniques of incremental view maintenance in databases.
We complement their argument with a concrete implementation of 
 incremental matching using semi-na\"ive evaluation~\citep{seminaive}.
Moreover, 
 the literature studied the ``proof'' problem on \egraphs~\citep{pp-congr,flatt2022small}:
 many domains require not only the optimized terms that are equivalent to the original terms,
 but also a proof why they are equivalent.
\revdel{The database counterpart of this problem, called provenance, is also studied by database researchers~\citep{prov-semiring,prov-souffle}.}
A future direction is to study proof generations for \egglog programs%
\revdel{~by incorporating algorithms developed in both fields}.

% On the other hand, some components of equality saturation are unique due to its workloads.
% Because rewrite rules do not usually saturate in equality saturation,
%  schedulers are used in practice to prioritize rewrites that are more likely to be fruitful~\citep{egg}.
% However, because programs in languages like Datalog are usually run to fixpoint,
%  the scheduling problem for these languages are less studied.

\paragraph{Datalog and Relational Databases}
% Our work is closely related to two lines of works on Datalog and relational databases.
Functional dependency repairs via lattice joins in \egglog is directly inspired by Flix~\citep{flix}.
Flix extends Datalog by allowing relations to be optionally annotated by lattices.
With this feature, 
 Flix is able to express many advanced program analyses algorithms efficiently.
% \Egglog slightly generalizes Flix by allowing multiple heads in the rules, 
% A difference between Flix and \egglog is that 
%  we use the standard equality joins 
%  while Flix uses meet-semantics for relational joins,
%  which can nonetheless be simulated in \egglog.
\egglog can simulate Flix programs by setting \merge to the lattice join operator.   
Flix can be regarded as among the works that 
 try to find a theoretical foundation for recursive aggregates~\citep{agg-semantics,datalog-subsumption,mono-agg,datalogo}.
A similar lattice-based approach to recursive aggregates is studied by $\text{Bloom}^L$~\citep{bloom-lattice}.
\revision{Other Datalog systems that support recursive aggregates include LogicBlox~\citep{logicblox} 
 and Rel~\citep{rel-refernce}.
}

Rewrite rules in \egglog generalize rules in Datalog 
 because the heads of rules can generate fresh ids.
This is called tuple-generating dependencies (TGDs) in the database literature.
Moreover, 
 while functional dependencies has the form $R(x_1,\ldots,x_i,\ldots,x_k),R(x_1,\ldots,x_i',\ldots x_k)\rightarrow x_i=x_i'$
 equality-generating dependencies (EGDs) generalize functional dependencies 
 by allowing equalities between different columns in different relations.
A family of algorithms called the chase can be used to reason about both TGDs and EGDs~\citep{chase-revisited,bench-chase}.
Moreover,
 the model semantics of the chase directly 
 gives a model semantics of a subset of \egglog where 
 union is the only \lstinline[language=egglog]{:merge} operation.
Compared to general TGDs and EGDs,
 \egglog imposes syntactic constraints over the programs 
 so that rules applications are deterministic and efficient.
$\text{Datalog}^{\pm}$~\citep{datalogpm} is 
 a family of extensions to Datalog based on TGDs and EGDs for ontological reasoning.

\revision{
Concurrent to our work and independent to the work on the chase, 
 \citet{bidlingmaier2023algebraic}
 formalizes Datalog with equality, which shares the same core idea to \egglog, 
 as relational Horn logic and partial Horn logic
 and studies its properties from a categorical point of view.
\citet{bidlingmaier2023evaluation} further describes an evaluation algorithm 
 for Datalog with equality similar to the chase.
Different from theirs, our work is motivated by practical applications 
 in program optimization and program analysis,
 and we focus on a simpler operational model of \egglog.
}

\revision{
While termination for Datalog with a variety of extensions is well studied, 
 the termination condition of \egglog is quite open.
In the future, 
 we hope to better understand the termination of \egglog by 
 further studying the connection between \egglog and the chase.
For example, 
 the database theory community has established many conditions 
 for chase termination (e.g., \citet{data-exchange,datalogpm,vadalog}),
 and one could potentially apply these results to \egglog
 by translating \egglog rewrite rules and functional dependencies into TGDs and EGDs.
On the other hand,
 nearly all instantiations of equality saturation in practice will diverge . 
Being a generalization of equality saturation, \egglog allows for divergence by design.}

\revision{
A key feature of \egglog is its efficient equational reasoning.
Although equational reasoning can be expressed in Datalog 
 with an explicit equivalence relation,
 doing so is very inefficient.
The \texttt{patched} baseline in \autoref{sec:points-to} 
 is a slightly more efficient encoding of equational reasoning in Datalog
 using \souffle extensions including choice domain and subsumptive rules.
We also attempted several other approaches to optimize equational reasoning 
 with existing features such as recursive aggregates and
 Constraint Handling Rule's simpagation rules \citep{chr}.
However, 
 we found none of these encodings 
 provide a natural abstraction nor competent performance.
}

% \yz{Other things to potentially talk about: 
%  functional programming in Datalog;
%  terms in Datalog;
%  formulog?}

\paragraph{Logic Programming and Automated Theorem Proving}
\revdel{
\egglog bears some similarity with logic programming languages like Prolog.
Prolog has unification variables similar to \egglog, except that they are backtrackable.
In general, Prolog allows disjunction and backtracking search
 at the cost of incomplete search strategy and potentially getting into infinite loops.
Moreover, Prolog programs are usually evaluated top-down,
 while Datalog and \egglog programs are evaluated bottom-up.
}

\revision{
\egglog bears some similarity with logic programming languages like Prolog.
Similar to unification variables in Prolog, 
 fresh ids in \egglog can represent unknown information 
 (see, e.g., \citet[Appendix~A.1]{egglog-preprint}), 
 and the congruence closure can be viewed 
 as a dual procedure to unification \citep{congr-duality}. 
In \citet[Appendix~A.3]{egglog-preprint}, we also show an 
 implementation of a Hindley-Milner type inference algorithm,
 of which the key construct is the unification mechanism implemented 
 as a few \egglog rules.

However, 
 several distinguishing features make \egglog highly efficient for its target domains, namely program analysis and optimization. 
For example,
 \egglog does not allow backtracking, 
 so its union-find data structure does not need to be backtrackable or persistent (unlike in Prolog or SMT solvers), 
 which makes it efficient for tasks that are monotone in nature 
 (e.g., equality saturation and pointer analysis).
Moreover, \egglog uses a bottom-up evaluation algorithm 
 more similar to Datalog than Prolog 
 (top-down backtracking search). 
One way of (partially) viewing \egglog is as a logic programming language that combines the bottom-up evaluation of Datalog
 and the unification mechanism of Prolog. The ``magic-set transformation'' is a closely related technique to simulate top-down evaluation 
 in a bottom-up language in a demand-driven fashion. 
We show in the full version \citep{egglog-preprint} several pearls that uses this idea to simulate top-down evaluations.
% We are continuing to explore to what extent this approach is expressive enough to implement key Prolog features.
On the other hand,
 Prolog has imperative features such as \texttt{cut}, which removes choice points.
 \Egglog does not have a direct analog of \texttt{cut}
 (because \egglog does not backtrack), 
 although \egglog has other imperative features borrowed from EqSat techniques
 that makes fine control of the execution such as rule scheduling.
}

SMT solvers are powerful tools for deciding combinations of logic theories and
 automatically proving theorems~\citep{z3,cvc4}.
Many rewrite rules in \egglog can be expressed as SMT axioms.
% For example, rule \lstinline[language=egglog]{(rewrite (Add a b) (Add b a))}
%  can be expressed as SMT axiom \lstinline{(forall ((x Expr) (y Expr)) (= (add x y) (add y x)))}.
In fact,
 SMT solvers support a richer language than \egglog with features like disjunction 
 and built-in theories like the theory of integer programming.
However, a key difference between \egglog and SMT solvers is that
 the output of \egglog is \emph{minimal} (or universal in database terminology).
% For instance, 
%  if we directly translate \autoref{fig:egglog-unification:b} into SMT axioms
%  and call an SMT solver,
%  the solver may return a trivial solution
%  by equating all expressions in the program (e.g., \lstinline{(Num 2)} and \lstinline{(Num 6)}) 
%  together and effectively having only one element in the universe.
% The solution by \egglog will instead make the least number of unifications,
%  which is more useful for many applications.
While it is possible to ``hack into'' an SMT solver to repurpose it 
 as an EqSat engine~\cite{flatt2022small}, 
 such techniques are arcane and not officially supported.
\eqsat and \egglog's native support for extraction
 makes them better suited for program optimization.
% \yz{I believe this is true, but I'm not 100\% sure about this. 
% Will a reviewer nitpick by giving a very smart encoding of eqsat in SMT?}
% \rw{Check with Oliver, he has abused SMT as egraphs.}
Recently, researchers have extended Datalog to dispatch more complex constraints 
 to be solved by an SMT solver~\cite{DBLP:journals/pacmpl/Bembenek0C20}.
This greatly extends the reach of Datalog, 
 allowing the user to specify constraints in a variety of theories and logics.
In \egglog we emphasize efficiency, 
 and chose more conservative extensions that can be implemented 
 by fast data structures like union-find.
